\section{Related Work}
\label{sec:related}

Detecting steppable surfaces in staircase point clouds touches several
threads of point-cloud processing. We organise the literature thematically
and weigh each line of work against the requirements of our problem:
fast, interpretable, and label-free segmentation of treads and risers.

\subsection{Geometric plane segmentation}
The dominant paradigm for extracting planar structure from unorganised
clouds is Random Sample Consensus (RANSAC). Li \emph{et al.}\
\cite{li2017improved} improve standard RANSAC for indoor and urban plane
segmentation by guiding sampling with Normal Distribution Transformation
cells, which mitigates the spurious planes that plain RANSAC fits across
disjoint surfaces. The same plane-fitting machinery appears in RGB-D
object recognition \cite{jalal2021rgb} and inside real-time planar LiDAR
SLAM \cite{zhou2021lsam}, where planes act as map primitives. Such
methods are efficient and interpretable. The difficulty is that they
treat planes as generic primitives with no orientation prior. On a
staircase the globally dominant plane is the oblique ramp spanning all
steps, so unconstrained fitting absorbs riser points into tread planes.
This observation motivates the orientation gate in our pipeline.

\subsection{Normal estimation and surface analysis}
A complementary route reasons about local surface geometry rather than
global primitives. Huang \emph{et al.}\ \cite{huang2009consolidation}
consolidate unorganised clouds and estimate reliable normals for surface
reconstruction, whereas Castillo \emph{et al.}\ \cite{castillo2013point}
recover normals through constrained nonlinear least squares for joint
segmentation and denoising. Per-point normals are attractive for stairs
because tread and riser differ sharply in orientation. On its own,
though, normal estimation produces no grouping into coherent surfaces and
remains sensitive to neighbourhood scale under noise. Our PCA-normal
baseline builds directly on this principle, classifying points by the
tilt of the smallest-eigenvector normal relative to the vertical axis.

\subsection{Clustering and region growing}
To recover surfaces rather than isolated points, clustering and
region-growing methods aggregate neighbours by geometric similarity.
Density-based clustering (DBSCAN) has been used to identify and localise
3D weld seams robustly in the presence of outliers
\cite{zhou2024identification}. Euclidean clustering supports object
detection from fused sparse clouds and imagery \cite{xu2021object}, and
supervoxel-based region growing yields compact, boundary-respecting
segments \cite{luo2021supervoxel}. Approaches of this kind naturally
reject outliers as noise and form one cluster per physical surface, which
suits the discrete tread and riser structure of stairs. What they do not
encode is the horizontal-versus-vertical distinction we require, so we
precede spatial clustering with an orientation split and confirm each
cluster by its slope.

\subsection{Ground and LiDAR segmentation}
Separating traversable ground from obstacles is a long-studied special
case of surface segmentation. The survey by Gomes \emph{et al.}\
\cite{gomes2023survey} catalogues ground-segmentation methods for
automotive LiDAR and highlights a recurring trade-off: geometric methods
are fast and transparent, while learned methods are accurate but opaque.
Stairs are a structured instance of this ground-versus-non-ground
problem, in which each tread is a small, elevated ground patch rather
than a single dominant plane.

\subsection{Learning-based segmentation}
Deep networks have reshaped point-cloud semantic segmentation since
PointNet \cite{qi2016pointnet} demonstrated direct learning on raw point
sets, and recent surveys document the breadth of this progress
\cite{sarker2024comprehensive}. Because dense labels are costly, weakly
supervised \cite{li2022hybridcr}, few-shot \cite{zhao2020few}, and
prototype-based few/zero-shot \cite{he2023prototype} formulations all aim
to reduce annotation demand. Accuracy is strong, but these models stay
data-hungry, depend on representative training distributions, and offer
limited interpretability and reproducibility for a safety-critical
sub-problem such as foothold detection. That cost in data and
transparency is precisely what motivates revisiting classical geometry
for stairs.

\subsection{Robotic stair and terrain perception}
In robotics, stair and terrain perception is tightly coupled to control.
Matsumura \emph{et al.}\ \cite{matsumura2022deep} detect stairs from 3D
point clouds with deep learning to prevent falls, and operator-assisted
stair traversal has been demonstrated in field systems
\cite{marion2017director}. Legged platforms lean on elevation maps and
adaptive planning \cite{belter2016adaptive}, egocentric vision
\cite{agarwal2022legged}, and stereo-based navigation over rough terrain
\cite{stelzer2012stereo}, while traversability itself can be learned from
simulation \cite{chavezgarcia2017learning}. Such systems confirm that
reliable steppable-surface detection is a prerequisite for safe motion.
They are typically evaluated end-to-end on bespoke platforms, however,
which leaves the underlying segmentation step without an isolated,
comparable benchmark.

\subsection{Gap addressed by this work}
Across these threads, geometric methods are fast and interpretable yet
rarely compared head-to-head on stairs, while learned methods require
labelled data that is unavailable for this specific task. No prior study
offers a controlled, fully labelled staircase benchmark on which
classical segmenters can be measured against a common ground truth. We
fill this gap with a parametric, fully annotated point-cloud generator
and an apples-to-apples comparison of four classical methods for
steppable-surface detection.
